

\section{Deep Learning in Robotics Fields}

%In this Section, we will review previous research on the application of deep learning in robotics and establish the context for our proposed concept of deep predictive learning.
In this section, we review previous research on the application of deep learning in robotics and establish the context for our proposed deep predictive learning concept.



\subsection{Recognition}

%The term "deep learning" was first introduced in the field of robotics at the IEEE/RSJ IROS 2017 conference, which was relatively late compared to other domains. Many early applications of deep learning in robotics focused on areas such as object detection in robot vision \cite{Ren2016pami,Redmon2016cvpr} and research in conversational speech recognition \cite{Hinton2013spm,Graves2014icml}. In traditional robotic intelligence design, the process typically involves the following steps: (1) recognizing the environment or objects, (2) correlating recognition results with pre-provided 2D maps or 3D CAD models, (3) understanding the current positions, object positions, and orientations, and (4) planning actions based on the situation. Deep learning significantly enhanced the first step, which is recognition capability. By learning labels like graspable regions that were previously challenging to model, it became possible to recognize a wide variety of general objects.
The term "deep learning" was introduced in robotics at the IEEE/RSJ IROS 2017 conference, which was relatively late compared with other domains. Several early applications of deep learning in robotics focused on object detection in robot vision~\cite{Ren2016pami,Redmon2016cvpr} and conversational speech recognition~\cite{Hinton2013spm,Graves2014icml}. In traditional robotic intelligence design, the process typically involves the following steps: (1) recognizing the environment or objects, (2) correlating the recognition results with pre-provided 2D maps or 3D CAD models, (3) understanding the current positions, object positions, and orientations, and (4) planning actions based on the situation. Deep learning significantly enhances the recognition capability in the first step. Learning labels, such as graspable regions that were previously challenging to model, made recognizing a wide variety of general objects possible.


%However, the burden of manual labeling by humans is substantial, and relying solely on image-based learning does not guarantee successful object grasping. Grasping performance depends on a combination of visual information, object attributes, and the robot's physical characteristics. Factors like the gripper's force, friction conditions, and object deformations must all be considered, and there are countless such elements to account for.
However, the burden of manual labeling by humans is substantial, and relying solely on image-based learning does not guarantee successful object grasping. The grasping performance depends on a combination of visual information, object attributes, and physical characteristics of the robot. Factors such as the gripper force, friction conditions, and object deformations must be considered, as well as countless elements.



\subsection{Robot Motion Learning}

%In response to the developments in perception systems mentioned above, there has been active research~\cite{Noda2014multimodal,Lee2019icra,lin2022tactilegym2} in End-to-End methods for predicting robot motion trajectories directly from multi-modal input sensor data (such as visual, auditory, tactile, etc.). Within the realm of learning-based approaches for robot motion, deep learning-based methods can be broadly categorized into two groups.
In response to the developments mentioned above in perception systems, active research has been conducted ~\cite{Noda2014multimodal,Lee2019icra,lin2022tactilegym2} on end-to-end methods for predicting robot motion trajectories directly from multimodal input sensor data (such as visual, auditory, tactile, etc.). Deep learning-based methods can be broadly categorized into two groups within learning-based approaches to robot motion.



\subsubsection{Reinforcement Learning}

%In deep reinforcement learning, the agent predicts rewards for combinations of high-dimensional inputs related to the environment and several executable action outputs by the robot. With each attempt by the agent, rewards are provided from the environment, and the model updates its parameters to acquire a policy that maximizes rewards~\cite{Levine2016jmlr,Levine2017ijrr,Haarnoja2017icml}. The application of deep reinforcement learning in games has been extremely successful~\cite{Mnih2013playing,Hessel2018rainbow}, and its application to robotics is a natural progression. Its greatest appeal lies in the significant reduction of the traditional process of designing models for the environment or objects, which designers had to provide. This has led to the proposal of various task-specific methods~\cite{Tsurumine2019ras,Mahler2019learning}. Levine and colleagues replaced previously required image feature engineering with deep learning and built an end-to-end model from images to actions~\cite{Levine2017ijrr}. On the other hand, in~\cite{lee2023pi,rlscale2023arxiv}, efforts were made to enhance the performance of general agents by training multiple tasks with a single model. However, achieving the optimal model for online learning in a real environment remains a challenge, as it requires a large number of trials and resources. There is also the approach of collecting a large amount of data with human resources~\cite{dasari2019robonet,rt12022arxiv}, but this approach has limitations in terms of data diversity.
In deep reinforcement learning, the agent predicts rewards for combinations of high-dimensional inputs related to the environment and several executable action outputs from the robot. With each attempt by the agent, the environment provides rewards, and the model updates its parameters to acquire a policy that maximizes the rewards~\cite{Levine2016jmlr,Levine2017ijrr,Haarnoja2017icml}. The application of deep reinforcement learning in games has been extremely successful~\cite{Mnih2013playing,Hessel2018rainbow}, and its application in robotics is a natural progression. Its greatest appeal lies in significantly reducing the traditional process of designing models for the environment or objects that designers must provide. This innovation has led to various task-specific methods~\cite{Tsurumine2019ras,Mahler2019learning}. Levine et al. replaced the previously required image feature engineering with deep learning and built an end-to-end model from images to actions~\cite{Levine2017ijrr}. However, in~\cite{lee2023pi,rlscale2023arxiv}, efforts were made to enhance the performance of general agents by training multiple tasks with a single model. However, achieving an optimal model for online learning in a real environment remains a challenge, as it requires many trials and resources. There is also an approach for collecting a large amount of data using human resources~\cite{dasari2019robonet,rt12022arxiv}, but this approach has limitations regarding data diversity.


%As a solution to the above issues, there is the Sim2Real approach that uses simulation~\cite{James2019cvpr,Chebotar2019icra}. By making simulation closer to reality~\cite{allevato2020tunenet} and transforming it~\cite{rao2020rl}, it supplements the lack of data. Domain randomization, especially regarding visual and physical parameters~\cite{Tobin2017}, is utilized in tasks such as multi-fingered hand manipulation~\cite{openai2018learning} and rough terrain walking of quadruped robots~\cite{lee2020learning}. In~\cite{lee2020learning}, teacher-student learning using environment information obtained only in simulation has been successful in real-world motion control, which is quite intriguing. Recent advancements in physical simulation technology can replicate scenes equivalent to the real world, but detailed modeling of contact states remains a challenge~\cite{drake,Narang2022rss}. Currently, this problem is being discussed as one of the subjects in the field of Surface Science.
The Sim2Real approach uses simulations to address these issues~\cite{James2019cvpr,Chebotar2019icra}. Making simulations closer to reality~\cite{allevato2020tunenet} and transforming them~\cite{rao2020rl} can supplement the lack of data. Domain randomization, especially regarding visual and physical parameters~\cite{Tobin2017}, is utilized in tasks such as multi-fingered hand manipulation~\cite{openai2018learning} and rough terrain walking of quadruped robots~\cite{lee2020learning}. In~\cite{lee2020learning}, teacher-student learning using environmental information obtained only in simulations was successful in real-world motion control, which is quite intriguing. Recent advancements in physical simulation technology can replicate scenes equivalent to the real world; however, detailed modeling of contact states remains a challenge~\cite{drake,Narang2022rss}. This problem is currently being discussed as a subject in Surface Science.


%Moreover, traditional reinforcement learning typically requires the pre-designing of reward functions, which poses challenges as the model's performance is heavily dependent on it. Inverse reinforcement learning introduces a process of estimating reward functions from human demonstrations~\cite{Peng2018,Tsurumine2019humanoids,Wang2021}, reducing the burden of manually designing rewards based on task domain knowledge. However, aspects such as policy behavior design still remain, and the flexibility of applied tasks on real machines is limited.
Moreover, traditional reinforcement learning typically requires pre-designing reward functions, which poses challenges because the model's performance depends heavily on it. Inverse reinforcement learning introduces a process for estimating the reward functions from human demonstrations~\cite{Peng2018,Tsurumine2019humanoids,Wang2021}, thereby reducing the burden of manually designing rewards based on task domain knowledge. However, aspects such as policy behavior design remain, and the flexibility of the tasks applied to real machines is limited.



\subsubsection{Behavioral Cloning}

%In recent years, there has been growing interest in offline learning to address the challenges of data scale arising from online policy updates in conventional reinforcement learning. Particularly, in Behavioral Cloning (or Imitation Learning), learning is conducted to replicate motion sequence data (joint angles, velocities, torques, etc.) obtained from robots~\cite{yu2018one,Zhao2023rss,jang2021bc,Lynch2022,guhur2023instruction,chi2023diffusionpolicy}. Zhao et al. used a low-cost, fine-grained teleoperation system and employed Transformers to enable learning across various tasks~\cite{Zhao2023rss,fu2024mobile}. In this framework, it is possible to directly acquire a motion generation model from motion experience if effective data for task execution is available. Additionally, the feature of not requiring complex reward design is highlighted. As a result, it has been applied to various high-difficulty tasks such as flexible object manipulation~\cite{Kawaharazuka2019icra} and human-cooperative tasks~\cite{Sasagawa2020ral}. 
Recently, there has been growing interest in offline learning to address the challenges of data scale arising from online policy updates in conventional reinforcement learning. In Behavioral Cloning (or Imitation Learning), learning is conducted to replicate the motion sequence data (joint angles, velocities, torques, etc.) obtained from robots~\cite{yu2018one,Zhao2023rss,jang2021bc,Lynch2022,guhur2023instruction,chi2023diffusionpolicy}. Zhao et al. used a low-cost fine-grained teleoperation system and employed transformers to enable learning across various tasks~\cite{Zhao2023rss,fu2024mobile}. In this framework, it is possible to acquire a motion generation model directly from the motion experience if effective data for task execution are available. In addition, the feature of not requiring a complex reward design is highlighted. Consequently, it has been applied to various high-difficulty tasks, such as flexible object manipulation~\cite{Kawaharazuka2019icra} and human-cooperative tasks~\cite{Sasagawa2020ral}. 


%As related technologies, we should also mention offline reinforcement learning. Offline reinforcement learning aims to acquire a good policy solely from a pre-collected dataset, and these studies can be interpreted as reward-conditioned Behavioral Cloning~\cite{actionablemodels2021arxiv,ma2022vip}. These approaches seek to improve learning efficiency by utilizing past experience data~\cite{kalashnikov2018qt} or by simultaneously collecting data from multiple robots~\cite{kalashnikov2021mt}. In particular, Chen et al. proposed a sequentially predictable offline reinforcement learning model using Transformers~\cite{chen2021decision}. They predict actions from a sequence of tokens representing rewards, states, and actions, achieving high performance in long-term tasks using the expected cumulative future rewards.
Regarding related technologies, we should also mention offline reinforcement learning. Offline reinforcement learning aims to acquire a good policy solely from a pre-collected dataset, and these studies can be interpreted as reward-conditioned Behavioral Cloning~\cite{actionablemodels2021arxiv,ma2022vip}. These approaches seek to improve the learning efficiency by utilizing experience data~\cite{kalashnikov2018qt} or by simultaneously collecting data from multiple robots~\cite{kalashnikov2021mt}. In particular, Chen et al. proposed a sequentially predictable offline reinforcement learning model using transformers~\cite{chen2021decision}. They predicted actions from a sequence of tokens representing rewards, states, and actions and achieved high performance in long-term tasks using expected cumulative future rewards.


%Many of the studies mentioned so far focus on modeling robot action sequences alone. In contrast, the "Predictive Learning" approach aims to perform self-supervised learning, including the external environment~\cite{lotter2016deep}. Lotter et al. proposed PredNet, which mimics predictive coding processes in the cerebral cortex. PredNet learns to minimize the error between predicted and actual values in video prediction tasks. These are implementations of the predictive coding theory~\cite{Rao1999} using deep learning models but do not consider interaction with the external environment in robotics.
Many of the studies mentioned thus far have focused solely on modeling robot action sequences. In contrast, the "Predictive Learning" approach aims to perform self-supervised learning, including the external environment~\cite{lotter2016deep}. Lotter et al. proposed PredNet, which mimics the predictive coding processes in the cerebral cortex. PredNet learns to minimize the error between the predicted and actual values in video prediction tasks. These are implementations of predictive coding theory~\cite{Rao1999} using deep learning models, but they do not consider interaction with the external environment in robotics.


%Our proposed concept, from the perspective of cognitive robotics~\cite{Asada2015}, incorporates predictive coding in sensorimotor dynamics. The crucial idea in our proposal is that the learning target of autonomous systems should not be a "sensor-to-motion map" but rather a "continuous loop of integrated sensorimotor dynamics." In systems where the environment and the body dynamically interact, the distinction between sensation and motion is always ambiguous. What exists there is a continuous flow where the current something leads to the future outcome. It is crucial to have an orientation to fill the error of the predictive flow that never completely fills the self and external. This kind of predictive learning system (forward model) and the corresponding motion generation system (inverse model) have been emphasized for a long time, but practical robots explicitly designed to work with this concept in the real world are still limited.
Our proposed concept, from the perspective of cognitive robotics~\cite{Asada2015}, incorporates predictive coding into sensorimotor dynamics. The crucial idea in our proposal is that the learning target of autonomous systems should not be a "sensor-to-motion map" but rather a "continuous loop of integrated sensorimotor dynamics." The distinction between sensation and motion is always ambiguous in systems where the environment and body interact dynamically. A continuous flow exists in which the current leads to future outcomes. It is crucial to have an orientation that fills the error of the predictive flow, which never completely fills the self and the external. This type of predictive learning system (forward model) and the corresponding motion generation system (inverse model) have been emphasized for a long time. However, practical robots designed to work with this concept in the real world remain limited.



